\documentclass[10pt]{beamer}
\usetheme{metropolis}
\usepackage{graphicx}\usepackage{listings}\usepackage{booktabs}\usepackage{xcolor}
\definecolor{codebg}{HTML}{F5F5F5}\definecolor{codegreen}{HTML}{2E7D32}\definecolor{codegray}{HTML}{757575}\definecolor{codeblue}{HTML}{1565C0}
\lstdefinestyle{rstyle}{language=R,backgroundcolor=\color{codebg},basicstyle=\ttfamily\scriptsize,keywordstyle=\color{codeblue}\bfseries,commentstyle=\color{codegray}\itshape,breaklines=true,frame=single,rulecolor=\color{codegray!30},showstringspaces=false,morekeywords={library,leaflet,addTiles,addProviderTiles,addCircleMarkers,addPolygons,addLegend,setView,addHeatmap,tmap,tmap_mode,plotly,plot_ly,ggplotly}}
\lstdefinestyle{pythonstyle}{language=Python,backgroundcolor=\color{codebg},basicstyle=\ttfamily\scriptsize,keywordstyle=\color{codeblue}\bfseries,commentstyle=\color{codegray}\itshape,breaklines=true,frame=single,rulecolor=\color{codegray!30},showstringspaces=false,morekeywords={import,from,as,folium,Map,Marker,CircleMarker,Choropleth,HeatMap,GeoJson,Popup,Tooltip,plugins,px,plotly}}
\title{Interactive Maps}
\subtitle{Workshop 5 --- Module 8}
\author{Prof.Asoc.\ Endri Raco}
\institute{Data Visualization for Data Scientists \\ UNYT Tirana}
\date{2025/26}
\begin{document}
\begin{frame}\titlepage\end{frame}
\begin{frame}{Module Roadmap}
\begin{enumerate}
\item Why interactive maps: pan, zoom, hover, click
\item The tile-map stack: base tiles + overlays
\item leaflet (R) and folium (Python): the core libraries
\item Markers, popups, tooltips, and layer control
\item Interactive choropleth and heatmaps
\item When interactive beats static (and vice versa)
\end{enumerate}
\begin{exampleblock}{Learning Outcome}
You will build interactive web maps with markers, popups, choropleth layers, and heatmaps using leaflet (R) and folium (Python), and export them as standalone HTML files.
\end{exampleblock}
\end{frame}
\section{Interactive Map Tools}
\begin{frame}{The Interactive Map Landscape}
\begin{center}\includegraphics[width=0.88\textwidth]{figures_w05m08/interactive_tools}\end{center}
\end{frame}
\begin{frame}{The Tile-Map Stack}
\centering
\begin{tabular}{llp{4.5cm}}
\toprule
\textbf{Layer} & \textbf{What} & \textbf{Example} \\
\midrule
Base tiles & Background map imagery & OpenStreetMap, CartoDB, Stamen \\
Vector overlay & Points, lines, polygons & Markers, GeoJSON boundaries \\
Raster overlay & Heatmap, choropleth fill & KDE surface, coloured regions \\
Controls & Legend, layer toggle, zoom & Layer switcher, scale bar \\
Interaction & Pan, zoom, click, hover & Popups, tooltips \\
\bottomrule
\end{tabular}
\end{frame}
\section{leaflet (R)}
\begin{frame}[fragile]{leaflet in R: Markers + Popups}
\begin{lstlisting}[style=rstyle]
library(leaflet)

leaflet(cities) |>
  addProviderTiles("CartoDB.Positron") |>      # light base tiles
  setView(lng = 19.82, lat = 41.33, zoom = 7) |>
  addCircleMarkers(
    lng = ~lon, lat = ~lat,
    radius = ~sqrt(population) / 100,            # proportional
    color = "#1565C0", fillOpacity = 0.5,
    popup = ~paste0("<b>", name, "</b><br>",      # HTML popup
                    "Pop: ", format(population, big.mark=",")),
    label = ~name) |>                             # hover tooltip
  addLegend(position = "bottomright",
    colors = "#1565C0", labels = "Cities", title = "Albania")

# Instant interactive from tmap
library(tmap); tmap_mode("view")
tm_shape(districts) + tm_fill("gdp") + tm_dots(cities)
\end{lstlisting}
\end{frame}
\begin{frame}[fragile]{leaflet: Choropleth + Heatmap}
\begin{lstlisting}[style=rstyle]
# Interactive choropleth
leaflet(districts) |>
  addProviderTiles("CartoDB.Positron") |>
  addPolygons(fillColor = ~pal(gdp),           # choropleth fill
    fillOpacity = 0.7, weight = 1, color = "white",
    highlight = highlightOptions(weight = 3),   # hover highlight
    popup = ~paste0(name, ": ", gdp, " EUR")) |>
  addLegend(pal = pal, values = ~gdp, title = "GDP/cap")

# Heatmap (requires leaflet.extras)
library(leaflet.extras)
leaflet(events) |>
  addProviderTiles("CartoDB.DarkMatter") |>
  addHeatmap(lng = ~lon, lat = ~lat, intensity = ~magnitude,
    blur = 20, radius = 15, max = 0.05)
\end{lstlisting}
\end{frame}
\section{folium (Python)}
\begin{frame}[fragile]{folium in Python: Markers + Popups}
\begin{lstlisting}[style=pythonstyle]
import folium

m = folium.Map(location=[41.33, 19.82], zoom_start=7,
    tiles="CartoDB positron")

for _, row in cities.iterrows():
    folium.CircleMarker(
        location=[row["lat"], row["lon"]],
        radius=row["population"] ** 0.5 / 100,
        color="#1565C0", fill=True, fill_opacity=0.5,
        popup=folium.Popup(
            f"<b>{row['name']}</b><br>Pop: {row['population']:,}"),
        tooltip=row["name"]
    ).add_to(m)

m.save("albania_cities.html")  # standalone HTML
\end{lstlisting}
\end{frame}
\begin{frame}[fragile]{folium: Choropleth + HeatMap}
\begin{lstlisting}[style=pythonstyle]
# Interactive choropleth
folium.Choropleth(
    geo_data="districts.geojson",
    data=districts_df,
    columns=["id", "gdp"],
    key_on="feature.properties.id",
    fill_color="YlOrRd",
    legend_name="GDP per capita (EUR)"
).add_to(m)

# Heatmap (folium.plugins)
from folium.plugins import HeatMap
heat_data = list(zip(events["lat"], events["lon"],
    events["magnitude"]))
HeatMap(heat_data, radius=15, blur=20,
    max_zoom=12).add_to(m)

m.save("albania_interactive.html")
\end{lstlisting}
\end{frame}
\section{Static vs Interactive}
\begin{frame}{When to Use Which}
\centering
\begin{tabular}{lll}
\toprule
\textbf{Context} & \textbf{Format} & \textbf{Why} \\
\midrule
PDF report / journal & Static & Print-ready, no tech dependency \\
Presentation slides & Static & Reproducible, no internet needed \\
Web dashboard & Interactive & Pan, zoom, hover = exploration \\
Exploratory analysis & Interactive & Drill into specific locations \\
Public communication & Interactive & Engagement, accessibility \\
Large $n$ ($>$50K) & Interactive + tiles & Lazy loading, WebGL render \\
\bottomrule
\end{tabular}
\vspace{6pt}
\begin{alertblock}{Pro-Tip}
Produce \textbf{both}: a static version for the PDF report and an interactive HTML alongside. In R, \texttt{tmap\_mode()} switches between the two with one command. In Python, produce \texttt{geopandas.plot()} for static and \texttt{folium} for interactive from the same data.
\end{alertblock}
\end{frame}
\section{Synthesis}
\begin{frame}{Key Takeaways}
\begin{enumerate}
\item Interactive maps add \textbf{pan, zoom, hover, click} --- exploration impossible in static maps.
\item The stack: \textbf{base tiles} (OSM, CartoDB) + \textbf{vector overlay} (markers, GeoJSON) + \textbf{controls} (legend, layer toggle).
\item \textbf{R}: \texttt{leaflet()} for full control, \texttt{tmap\_mode("view")} for instant switch from static.
\item \textbf{Python}: \texttt{folium.Map()} + \texttt{.add\_to()} chain. \texttt{folium.plugins.HeatMap()} for density.
\item Export as \textbf{standalone HTML} (\texttt{m.save("map.html")}): no server needed, shareable by email.
\item Produce \textbf{both static and interactive} versions for every spatial analysis.
\end{enumerate}
\end{frame}
\begin{frame}{Essential Readings}
\begin{itemize}
\item leaflet for R: \texttt{https://rstudio.github.io/leaflet/}
\item folium: \texttt{https://python-visualization.github.io/folium/}
\item Lovelace, R. et al. (2019). \emph{Geocomputation with R}, Chapter 8.
\end{itemize}
\vspace{6pt}
\textbf{Next Module}: Snow's Cholera Map --- Saving Lives with Data (W05-M09)
\end{frame}
\end{document}
